\paragraph{Opportunity accounting.} The model's opportunities are patient-zone
boundary events: entry, where WHO Moment 1 falls, and exit, where Moments 4 and 5
fall; Moments 2 and 3 inside the zone are not represented \citep{who2009}. WHO
direct observation counts coinciding indications as one opportunity requiring one
action; in particular a move from patient A directly to patient B, which the model
scores as two boundary events (exit A, then entry B, with no contact between), is
one WHO opportunity, met if either action occurs \citep[Sections II.5.6 and
III.5]{who2009}. The WHO \emph{Guidelines} say when this merger applies: Moments 4
and 1 typically coincide ``when moving directly from one patient to another
without touching any surface outside the corresponding patient zones''
\citep[Part I, Section 21.4.2, p.~102]{whoguidelines2009}; under their
definitions of the two moments, touching a health-care-area surface between
patients makes exit and entry separate opportunities.
The model has no contactable surfaces outside care zones, so every model move is
direct. On a real ward, audited compliance therefore mixes one-opportunity direct
moves with two-opportunity indirect moves (besides Moments 2 and 3), in a
proportion the model does not represent, and neither the per-boundary nor the
transition definition used here maps onto an audited figure without the share of
direct moves. Its miss probability is $(1-h)\{\rho+(1-\rho)(1-h)\}$, which
equals $q_2$ when $e=1$ (Table~\ref{tab:who}). A WHO-style audit of direct
transitions therefore estimates the behavioural part of $q_2$, and boundary-scored
audits (as in some electronic systems) estimate $h$. Electronic systems differ in
what they score: a system that credits one dispense to an exit and the next entry
within a time window moves its score toward transition units; independent
misclassification of events would attenuate an estimated $\rho$; shared bays
may be harder to instrument than single rooms; and group-level systems give $h$
without per-worker sequences. $\rho$ should therefore be estimated from raw
per-worker event streams, not from vendor compliance scores. Under independence a
transition compliance $c$ corresponds to per-boundary $h=1-\sqrt{1-c}$, so
transition compliance of 90\% and 95\% corresponds to $h\approx.684$ and .776. If a
policy were defined as transition compliance rising from 90\% to 95\% with perfect
clearing, two-opportunity survival would fall from .10 to .05 and $F_2=50\%$ for
every $\rho$. In the full-ICU coefficient (prespecified extension,
Section~\ref{sec:s-coef-extension}, Table~\ref{tab:extension-transition}) this
policy gave 52.36\% ($\rho=0$) and 52.48\% ($\rho=.5$) with perfect clearing, a
clustering contrast of $+0.12$ percentage points, but 34.18\% and 41.62\% at
$e=.8$ ($+7.45$ points); the per-boundary .90$\to$.95 policy gives $\mn21.5$
and $+9.0$ points on the same backgrounds. The two policies are different
interventions, with different per-boundary completion levels, so only the
comparison across $\rho$ within one compliance definition is like-for-like;
which compliance definition sets the policy target changes the forecast. Both
definitions are real measurements: boundary-scored electronic monitoring of
individuals records per-boundary events, whereas WHO direct observation scores a
direct patient-to-patient move as one opportunity. Their
correspondence here rests on model assumptions (the stationary Markov law and
one exit and one entry draw per direct move), and WHO-style transition audits
alone determine neither $h$ nor $e$. Elsewhere the $\rho$-dependence is
relative to a policy that raises per-boundary completion with $\rho$ fixed. Holding instead the
miss-to-complete probability or the mean miss-run length fixed would also change
$F$ for the same $h_0\to h_1$. Runs of consecutive misses occur under independence
too, with geometric lengths; models in which risk grows with consecutive misses
can assume independent compliance with geometric run lengths \citep{chan2026}.
Runs alone are therefore not evidence of serial dependence; the relevant
statistic is the lag-one (and, to check the Markov law, lag-two)
autocorrelation of unselected completion records.

